\section{Second Pass: Analysis as Quantifier Discipline}
\chapterline

In Chapter~34, we used mathematical induction to establish statements for all natural numbers: verifying a base case and stepping from $n$ to $n+1$. On the real number line, however, there is no ``next'' real number. Between any two distinct real numbers lie infinitely many others, so induction along a discrete succession cannot directly establish properties across a continuum.

To analyze limits, continuity, and change without contradiction, mathematicians formulated precise conditions using quantifiers. In elementary calculus, you compute limits using algebraic recipes: factor the numerator, cancel $(x-a)$, plug in numbers, or apply L'H\^opital's rule. 

In real analysis, definitions are stated directly using quantifiers:
\[
\forall \eps > 0, \; \exists \delta > 0, \; \forall x \in D, \quad |x - c| < \delta \implies |f(x) - f(c)| < \eps.
\]
Analysis formalizes intuitive notions like ``approaching,'' ``infinitely close,'' and ``continuous'' into precise logical statements governed by quantifier order.

Figure~\ref{35:fig-epsilon-delta-game} illustrates this adversarial contract as a two-player game. An adversary challenges you with a vertical tolerance window $L \pm \eps$. To prove continuity, you must respond with a horizontal window $c \pm \delta$ that traps the entire curve inside the adversary's corridor.

\begin{figure}[htbp]
\centering
\begin{tikzpicture}[scale=1.15, >=Latex]
  % Epsilon tolerance corridor
  \fill[targetgreen!12] (0, 1.2) rectangle (4.6, 2.8);
  \draw[dashed, targetgreen!80!black, thick] (0, 1.2) -- (4.6, 1.2);
  \draw[dashed, targetgreen!80!black, thick] (0, 2.8) -- (4.6, 2.8);

  % Delta control strip
  \fill[sourceblue!8] (1.5, 0) rectangle (3.5, 3.4);
  \draw[dashed, sourceblue!70!black, thick] (1.5, 0) -- (1.5, 3.4);
  \draw[dashed, sourceblue!70!black, thick] (3.5, 0) -- (3.5, 3.4);

  % Intersection target box
  \draw[very thick, draw=targetgreen!70!black, rounded corners=2pt] (1.5, 1.2) rectangle (3.5, 2.8);

  % Axes
  \draw[->, thick, gray!70] (-0.2, 0) -- (5.0, 0) node[right, font=\small] {$x$};
  \draw[->, thick, gray!70] (0, -0.2) -- (0, 3.7) node[above, font=\small] {$y$};

  % Function curve: smooth curve through (2.5, 2.0)
  \draw[very thick, structurepurple!90!black, domain=0.8:4.2, samples=50] 
    plot (\x, {2.0 + 0.6*(\x - 2.5) - 0.05*(\x - 2.5)^3});
  \node[font=\footnotesize\bfseries, text=structurepurple!90!black] at (4.3, 3.1) {$y = f(x)$};

  % Center point (c, L)
  \fill[structurepurple!90!black] (2.5, 2.0) circle (2.5pt);
  \draw[dotted, gray!60, thick] (2.5, 0) -- (2.5, 2.0) -- (0, 2.0);

  % Y-axis ticks and annotations (Adversary)
  \node[left=3pt, font=\scriptsize\bfseries, text=targetgreen!90!black] at (0, 2.8) {$L + \eps$};
  \node[left=3pt, font=\scriptsize\bfseries] at (0, 2.0) {$L = f(c)$};
  \node[left=3pt, font=\scriptsize\bfseries, text=targetgreen!90!black] at (0, 1.2) {$L - \eps$};
  
  % Adversary tolerance bracket on right side
  \draw[<->, thick, targetgreen!80!black] (4.8, 1.2) -- node[right=2pt, font=\tiny\bfseries, text=targetgreen!90!black, align=center] {Adversary\\tolerance\\$\pm \eps$} (4.8, 2.8);

  % X-axis ticks and annotations (Defender)
  \node[below=3pt, font=\scriptsize\bfseries, text=sourceblue!90!black] at (1.5, 0) {$c - \delta$};
  \node[below=3pt, font=\scriptsize\bfseries] at (2.5, 0) {$c$};
  \node[below=3pt, font=\scriptsize\bfseries, text=sourceblue!90!black] at (3.5, 0) {$c + \delta$};
  \draw[<->, thick, sourceblue!80!black] (1.5, -0.6) -- node[below=1pt, font=\tiny\bfseries, align=center] {Defender responds with window $\delta(\eps)$} (3.5, -0.6);
\end{tikzpicture}
\caption{The $\eps$-$\delta$ quantifier game. An adversary specifies an arbitrary vertical tolerance $\eps > 0$ around the limit $L = f(c)$ (green corridor). The prover wins the round by providing a horizontal control window $\delta > 0$ centered at $c$ (blue strip) such that every point on the curve within $(c - \delta, c + \delta)$ is trapped entirely inside the $(L - \eps, L + \eps)$ corridor.}
\label{35:fig-epsilon-delta-game}
\end{figure}

\begin{historybox}
For nearly two centuries after Isaac Newton and Gottfried Wilhelm Leibniz invented calculus in the late 1600s, mathematics operated on informal foundations. Calculations relied on ``infinitesimals''---quantities treated as smaller than any positive real number, yet nonzero when dividing. In 1734, the philosopher Bishop George Berkeley published \emph{The Analyst}, pointing out the logical contradiction of treating these increments as nonzero when dividing a difference quotient and as zero immediately afterward when discarding remainder terms.

The rescue operation culminated in nineteenth-century Germany. In 1872, Richard Dedekind published \emph{Stetigkeit und irrationale Zahlen} (\emph{Continuity and Irrational Numbers}). Dedekind defined real numbers without geometry or motion by partitioning the rationals into ``cuts'' ($A | B$)---splitting the rational line into left and right sets with no gap, thereby defining each irrational number as the unique divide between them. Simultaneously, Karl Weierstrass banished infinitesimals through the ``arithmetization of analysis,'' replacing fuzzy temporal phrases like ``$x$ approaches $c$'' with static, adversarial $\eps$-$\delta$ quantifier contracts: an adversary specifies a desired vertical tolerance $\eps$, and the prover must supply a horizontal control window $\delta$. Shortly thereafter, Eduard Heine (1872) and \'Emile Borel (1895) formulated what is now called the Heine--Borel theorem: on a closed and bounded interval, any collection of open intervals that covers the set can be thinned down to a finite subcollection, turning infinite challenges into finite checks.

Did infinitesimals die forever? In 1960, the logician Abraham Robinson used mathematical logic (model theory) to create \emph{non-standard analysis}, showing that infinitesimals can be made completely rigorous via the hyperreal numbers. Meanwhile, in modern \emph{computable analysis}, fundamental questions remain active: deciding whether an algorithmically defined real number equals zero is provably undecidable, keeping the computational nature of the continuum an open frontier.
\end{historybox}

\subsection{The Quantifier Game: Pointwise versus Uniform}

To see how drastically quantifier order dictates mathematical reality, compare two definitions of continuity for a function $f: D \to \R$ on a domain $D \subseteq \R$. Continuity means that sufficiently small movements in input guarantee that the output stays within any required tolerance:

\begin{itemize}
  \item \textbf{Pointwise Continuity on $D$}:
  \[
  \forall x \in D, \; \forall \eps > 0, \; \exists \delta > 0, \; \forall y \in D, \quad |x - y| < \delta \implies |f(x) - f(y)| < \eps.
  \]
  \item \textbf{Uniform Continuity on $D$}:
  \[
  \forall \eps > 0, \; \exists \delta > 0, \; \forall x \in D, \; \forall y \in D, \quad |x - y| < \delta \implies |f(x) - f(y)| < \eps.
  \]
\end{itemize}

Notice the microscopic shift: in pointwise continuity, the phrase ``$\forall x \in D$'' sits \emph{before} ``$\exists \delta > 0$''. In uniform continuity, ``$\exists \delta > 0$'' sits \emph{before} ``$\forall x \in D$''.

\begin{intuitionbox}
\textbf{The Order of Quantifiers:}
\begin{itemize}
  \item In \textbf{pointwise continuity}, the adversary provides a specific point $x$ and a tolerance $\eps$. The choice $\delta(x, \eps)$ may depend on the point $x$. If the slope grows steep near a boundary, $\delta$ can be chosen arbitrarily small near that location.
  \item In \textbf{uniform continuity}, the adversary provides only $\eps$. You must produce a single $\delta(\eps)$ that works across the entire domain $D$. If the slope grows arbitrarily steep anywhere in $D$, no single positive $\delta$ suffices globally, and uniform continuity fails.
\end{itemize}
\end{intuitionbox}

To see this difference in concrete arithmetic, consider $f(x) = 1/x$ on the open interval $(0, 1)$. The function is continuous at every single point $x_0 \in (0, 1)$. For any input $x \in (0, 1)$, the vertical difference is
\[
|f(x) - f(x_0)| = \left|\frac{1}{x} - \frac{1}{x_0}\right| = \frac{|x - x_0|}{x x_0}.
\]
If we keep $x$ close to $x_0$, say $|x - x_0| < x_0/2$, then $x > x_0/2$, which makes the denominator $x x_0 > x_0^2/2$. Thus
\[
|f(x) - f(x_0)| < \frac{2|x - x_0|}{x_0^2}.
\]
To force this change to be strictly smaller than $\eps$, we set $\delta = \min\{x_0/2, \, \eps x_0^2 / 2\}$. That proves pointwise continuity: for any chosen $x_0$, a valid $\delta > 0$ exists.

Yet $f(x) = 1/x$ is \emph{not} uniformly continuous on $(0, 1)$. Notice that our formula for $\delta$ shrinks proportionally to $x_0^2$. As $x_0 \to 0$, the curve becomes infinitely steep, and the required window collapses to zero. If an adversary challenges us with $\eps = 1$, no single universal $\delta > 0$ can work. For any candidate $\delta > 0$, choose an integer $n \in \N$ such that $1/n < \delta$. Now test the pair $x = 1/n$ and $y = 1/(2n)$ in $(0, 1)$. Their horizontal separation is
\[
|x - y| = \frac{1}{2n} < \frac{1}{n} < \delta,
\]
yet their vertical separation is
\[
|f(x) - f(y)| = |n - 2n| = n \ge 1 = \eps.
\]
The single $\delta$ fails. Hence $f(x) = 1/x$ is pointwise continuous on $(0, 1)$, but not uniformly continuous.

Figure~\ref{35:fig-pointwise-vs-uniform-hyperbola} plots this phenomenon directly on the graph of $f(x) = 1/x$. In the flat rightward region, a modest horizontal window $\delta_1$ easily controls a given tolerance $\eps$. But as $x \to 0$, the steepening gradient forces the required window $\delta_2$ to shrink toward zero. No single uniform $\delta$ can serve every point in $(0, 1)$ simultaneously.

\begin{figure}[htbp]
\centering
\begin{tikzpicture}[xscale=8.5, yscale=1.1, >=Latex]
  % Axes
  \draw[->, thick, gray!70] (-0.04, 0) -- (1.15, 0) node[right, font=\small] {$x$};
  \draw[->, thick, gray!70] (0, -0.2) -- (0, 5.4) node[above, font=\small] {$y$};
  
  \node[below left=2pt, font=\scriptsize] at (0, 0) {$0$};
  \draw[thick, gray!70] (1.0, 0) -- (1.0, -0.08) node[below=3pt, font=\scriptsize] {$1$};

  % Function curve y = 1/x
  \draw[very thick, structurepurple!90!black, domain=0.20:1.0, samples=60] plot (\x, {1/\x});
  \node[font=\footnotesize\bfseries, text=structurepurple!90!black, anchor=south west] at (0.20, 5.0) {$f(x) = 1/x$};

  % POINT 1: Gentle region at x1 = 0.6 (y1 = 1.67)
  \fill[sourceblue!15] (0.45, 1.07) rectangle (0.80, 2.27);
  \draw[thick, sourceblue!70!black, dashed] (0, 1.07) -- (0.80, 1.07);
  \draw[thick, sourceblue!70!black, dashed] (0, 2.27) -- (0.80, 2.27);
  \draw[thick, sourceblue!70!black, dotted] (0.45, 0) -- (0.45, 2.27);
  \draw[thick, sourceblue!70!black, dotted] (0.80, 0) -- (0.80, 2.27);
  
  \node[circle, fill=sourceblue!90!black, inner sep=2pt] at (0.60, 1.67) {};
  \node[below=3pt, font=\scriptsize\bfseries, text=sourceblue!90!black] at (0.60, 0) {$x_1$};
  \draw[<->, thick, sourceblue!90!black] (0.45, 0.3) -- node[above=1pt, font=\tiny\bfseries] {$\delta_1$ (wide)} (0.80, 0.3);

  % POINT 2: Steep region at x2 = 0.26 (y2 = 3.85)
  \fill[quotientorange!15] (0.225, 3.25) rectangle (0.308, 4.45);
  \draw[thick, quotientorange!80!black, dashed] (0, 3.25) -- (0.31, 3.25);
  \draw[thick, quotientorange!80!black, dashed] (0, 4.45) -- (0.31, 4.45);
  \draw[thick, quotientorange!80!black, dotted] (0.225, 0) -- (0.225, 4.45);
  \draw[thick, quotientorange!80!black, dotted] (0.308, 0) -- (0.308, 4.45);
  
  \node[circle, fill=quotientorange!90!black, inner sep=2pt] at (0.26, 3.85) {};
  \node[below=3pt, font=\scriptsize\bfseries, text=quotientorange!90!black] at (0.26, 0) {$x_2$};
  \draw[<->, thick, quotientorange!90!black] (0.225, 0.3) -- (0.308, 0.3);
  \node[above=2pt, font=\tiny\bfseries, text=quotientorange!90!black] at (0.26, 0.35) {$\delta_2$ (narrow)};

  % Tolerance brackets on y-axis
  \draw[<->, thick, sourceblue!90!black] (-0.06, 1.07) -- node[left=1pt, font=\tiny\bfseries] {$2\eps$} (-0.06, 2.27);
  \draw[<->, thick, quotientorange!90!black] (-0.06, 3.25) -- node[left=1pt, font=\tiny\bfseries] {$2\eps$} (-0.06, 4.45);

  % Explanatory banner placed cleanly in upper right
  \node[draw=gray!30, fill=gray!4, rounded corners=3pt, font=\scriptsize, align=center, text=gray!85!black] at (0.74, 4.4) {
    \textbf{Same tolerance $\eps$, different locations:}\\
    At $x_1$, gentle slope gives generous $\delta_1$.\\
    Near $0$, steep slope forces collapsing $\delta_2 \to 0$.\\
    No uniform $\delta$ exists on $(0, 1)$.
  };
\end{tikzpicture}
\caption{Why uniform continuity fails on $f(x) = 1/x$ across $(0, 1)$. For a fixed vertical tolerance $\eps$, the required control window $\delta(x, \eps) \approx \eps x^2$ shrinks quadratically as $x$ approaches $0$. At $x_1 = 0.6$, a wide window $\delta_1$ suffices; near $x_2 = 0.26$, the increasing slope compresses the allowable window to a sliver $\delta_2$. Because $\delta \to 0$ as $x \to 0$, no single positive window $\delta$ can serve the entire domain simultaneously.}
\label{35:fig-pointwise-vs-uniform-hyperbola}
\end{figure}

\subsection{Cauchy's Sum Theorem and the Shock of ``Monster'' Curves}

The struggle to formalize analysis produced some of the most famous false theorems and shocks in mathematical history:

\begin{itemize}
  \item \textbf{Cauchy's False Sum Theorem (1821)}: Augustin-Louis Cauchy, one of the greatest mathematicians in history, published a ``proof'' in his 1821 \emph{Cours d'Analyse} asserting that the limit of any convergent sequence of continuous functions is always continuous. Five years later in 1826, Niels Henrik Abel exhibited a glaring counterexample: the Fourier series $\sum_{n=1}^\infty \frac{\sin(nx)}{n}$---an infinite sum of sines and cosines used to model periodic waveforms---converges to a discontinuous sawtooth wave with periodic vertical jumps.
  
  The modern textbook counterexample is even simpler: consider $f_n(x) = x^n$ on the closed interval $[0, 1]$. Each function $f_n$ is a smooth, continuous curve passing from $(0, 0)$ to $(1, 1)$. For any fixed $x \in [0, 1)$, $\lim_{n \to \infty} x^n = 0$. But at $x = 1$, $1^n = 1 \to 1$. The pointwise limit function is
  \[
  f(x) = \begin{cases} 0 & \text{if } 0 \le x < 1, \\ 1 & \text{if } x = 1, \end{cases}
  \]
  which has a jump discontinuity at $x = 1$. A sequence of continuous curves has converged pointwise to a discontinuous limit function.
  
  For decades, mathematicians argued over what Cauchy had overlooked. In 1847, Philipp Ludwig von Seidel and George Gabriel Stokes isolated the missing concept, later perfected by Weierstrass: \emph{uniform convergence}. Cauchy had assumed pointwise convergence:
  \[
  \forall x \in D, \; \forall \eps > 0, \; \exists N \in \N, \; \forall n \ge N, \quad |f_n(x) - f(x)| < \eps,
  \]
  where $N$ can depend on both $\eps$ and the location $x$. But preserving continuity requires uniform convergence:
  \[
  \forall \eps > 0, \; \exists N \in \N, \; \forall x \in D, \; \forall n \ge N, \quad |f_n(x) - f(x)| < \eps,
  \]
  where a single index $N$ traps the entire curve inside an $\eps$-corridor around $f(x)$ simultaneously. Swapping two adjacent quantifiers had created a fifty-year controversy.

  \begin{figure}[htbp]
  \centering
  \begin{tikzpicture}[xscale=7.0, yscale=4.5, >=Stealth]
    % Tolerance corridor
    \fill[targetgreen!12] (0, 0) rectangle (1.0, 0.25);
    \draw[dashed, targetgreen!80!black, thick] (0, 0.25) -- (1.0, 0.25) node[right=2pt, font=\scriptsize\bfseries, text=targetgreen!90!black] {$+\eps = 0.25$};

    % Axes
    \draw[->, thick, gray!70] (-0.05, 0) -- (1.18, 0) node[right, text=black, font=\small] {$x$};
    \draw[->, thick, gray!70] (0, -0.05) -- (0, 1.25) node[above, text=black, font=\small] {$y$};

    \node[below left=2pt, font=\scriptsize] at (0, 0) {$0$};
    \draw[thick, gray!70] (1.0, 0) -- (1.0, -0.02) node[below=3pt, font=\scriptsize, text=black] {$1$};
    \draw[thick, gray!70] (-0.015, 1.0) -- (0, 1.0) node[left=3pt, font=\scriptsize, text=black] {$1$};

    % Function curves
    \draw[thick, sourceblue!50, domain=0:1, samples=40] plot (\x, {\x});
    \node[font=\scriptsize\bfseries, text=sourceblue!70!black] at (0.38, 0.48) {$n=1$};

    \draw[thick, sourceblue!70, domain=0:1, samples=40] plot (\x, {\x^2});
    \node[font=\scriptsize\bfseries, text=sourceblue!80!black] at (0.58, 0.44) {$n=2$};

    \draw[thick, sourceblue!90, domain=0:1, samples=50] plot (\x, {\x^5});
    \node[font=\scriptsize\bfseries, text=sourceblue!90!black] at (0.75, 0.40) {$n=5$};

    \draw[very thick, structurepurple, domain=0:1, samples=70] plot (\x, {\x^20});
    \node[font=\scriptsize\bfseries, text=structurepurple, anchor=south east] at (0.88, 0.12) {$n=20$};

    % Callout for corridor escape
    \draw[->, thick, trapred] (0.55, 0.70) to[bend left=12] (0.86, 0.32);
    \node[above, font=\scriptsize, text=trapred!90!black, align=center] at (0.55, 0.70) {Escapes $\eps$-corridor\\for every $n$};

    % Pointwise limit
    \draw[line width=2.5pt, targetgreen!90!black] (0, 0) -- (0.97, 0);
    \draw[thick, targetgreen!90!black, fill=white] (1.0, 0) circle (0.022);
    \fill[targetgreen!90!black] (1.0, 1.0) circle (0.022);
    \node[right=3pt, font=\scriptsize\bfseries, text=targetgreen!90!black] at (1.0, 1.0) {$(1, 1)$};
    \node[below=14pt, font=\scriptsize\bfseries, text=targetgreen!90!black] at (0.5, 0) {Pointwise limit $f(x) = 0$ on $[0, 1)$};
  \end{tikzpicture}
  \caption{Why Cauchy's Sum Theorem failed: the breakdown of uniform convergence for $f_n(x) = x^n$ on $[0, 1]$. Each curve $f_n(x)$ is continuous and unbroken from $(0, 0)$ to $(1, 1)$. For every fixed $x \in [0, 1)$, $x^n \to 0$ as $n \to \infty$; at $x = 1$, $1^n = 1$. The pointwise limit function (green) has a jump discontinuity at $x = 1$. If we choose a tolerance band $\eps = 0.25$ around the limit $0$, every single curve $x^n$ must eventually cross the boundary to reach $(1, 1)$, no matter how large $n$ is. Pointwise convergence holds at every individual coordinate, but uniform convergence fails because no single index $N$ keeps the entire curve inside the $\eps$-corridor simultaneously.}\label{35:fig-cauchy-counterexample}
  \end{figure}

  \item \textbf{Weierstrass's Nowhere-Differentiable Function (1872)}: Many nineteenth-century mathematicians believed that any continuous curve, except perhaps at isolated sharp corners, must possess tangent lines. In 1872, Weierstrass presented
  \[
  f(x) = \sum_{n=0}^\infty a^n \cos(b^n \pi x),
  \]
  where $0 < a < 1$, $b$ is an odd positive integer, and $ab > 1 + \frac{3}{2}\pi$. This function is continuous everywhere on $\R$, yet differentiable \emph{nowhere}---meaning it has no tangent line and no well-defined rate of change at any point. The curve displays self-similar oscillations under every magnification. Charles Hermite wrote to Thomas Stieltjes: ``I turn away with fright and horror from this lamentable evil of functions which have no derivatives!'' Only strict quantifier discipline could prove that such a function satisfies the definition of continuity while failing the definition of differentiability at every point.
  \item \textbf{The Reversal of Implication in Limit Scratchwork}: A persistent false start among students is confusing backwards exploration with forward proof. Manipulating the desired conclusion $|f(x) - L| < \eps$ backwards until you isolate $|x - c| < \delta$ merely shows that the conclusion implies the precondition ($Q \implies P$). Deduction must run strictly in the opposite direction ($P \implies Q$): you must demonstrate that whenever $0 < |x - c| < \delta$, the desired bound $|f(x) - L| < \eps$ is guaranteed to hold.
\end{itemize}

\subsection{The Nested Interval Theorem and Completeness}

Why are the real numbers $\R$ the standard setting for calculus, whereas the rational numbers $\Q$ are insufficient? The answer lies in \textbf{completeness}. Completeness means that the continuum has no gaps: every Cauchy sequence of real numbers converges to a limit that is itself a real number.

One of the most geometric expressions of completeness is the \textbf{Nested Interval Theorem}.

\begin{theorem}[Nested Interval Theorem]
Let $(I_n)_{n=1}^\infty$ be a sequence of closed, bounded intervals $I_n = [a_n, b_n]$ in $\R$ such that:
\begin{enumerate}
  \item The intervals are nested: $I_1 \supseteq I_2 \supseteq I_3 \supseteq \cdots$
  \item Their lengths shrink to zero: $\lim_{n \to \infty} (b_n - a_n) = 0$.
\end{enumerate}
Then their intersection contains exactly one point:
\[
\bigcap_{n=1}^\infty I_n = \{c\}, \qquad \text{where } c = \sup_{n} a_n = \inf_{n} b_n.
\]
\end{theorem}

\begin{proof}
Because the intervals are nested ($I_1 \supseteq I_2 \supseteq I_3 \supseteq \cdots$), their left endpoints $a_n$ move rightward (non-decreasing) while their right endpoints $b_n$ move leftward (non-increasing):
\[
a_1 \le a_2 \le a_3 \le \cdots \le b_3 \le b_2 \le b_1.
\]
In particular, every left endpoint $a_n$ is less than or equal to every right endpoint $b_m$. (If $n \le m$, then $a_n \le a_m \le b_m$; if $n > m$, then $a_n \le b_n \le b_m$.)

Therefore, the set of left endpoints $A = \{a_n : n \in \N\}$ is non-empty and bounded above by every $b_m$. By the completeness of the real numbers (the Least Upper Bound Property), $A$ has a supremum in $\R$:
\[
c = \sup_{n} a_n.
\]
By definition of supremum, $a_n \le c$ for all $n \in \N$. Also, because every $b_m$ is an upper bound for $A$, the least upper bound cannot exceed any $b_m$, so $c \le b_m$ for all $m \in \N$. Combining these inequalities gives
\[
a_n \le c \le b_n \quad \text{for all } n \in \N,
\]
which proves that $c \in I_n$ for every $n$, so $c \in \bigcap_{n=1}^\infty I_n$.

To see that $c$ is the \emph{only} point in the intersection, suppose $c$ and $d$ both belong to $\bigcap_{n=1}^\infty I_n$ with $c < d$. Then for every $n$, both points lie inside $[a_n, b_n]$, which forces the interval width to satisfy
\[
b_n - a_n \ge d - c > 0.
\]
Taking the limit as $n \to \infty$ contradicts hypothesis~(2) that $\lim_{n \to \infty} (b_n - a_n) = 0$. Hence the intersection contains exactly one point.
\end{proof}

\begin{figure}[htbp]
\centering
\begin{tikzpicture}[scale=1.15, >=Stealth]
  \tikzset{
    intervalbar/.style={very thick, rounded corners=2pt},
    endpoint/.style={circle, fill, inner sep=2pt},
    guide/.style={dashed, gray!60, thick}
  }

  % Interval 1
  \node[left=0.2cm, font=\small\bfseries, text=sourceblue!80!black] at (0.5, 3.4) {$I_1 = [a_1, b_1]$};
  \draw[intervalbar, draw=sourceblue!80!black, fill=sourceblue!15] (0.5, 3.2) rectangle (9.5, 3.6);
  \node[endpoint, fill=sourceblue!90!black, label={above left:\footnotesize $a_1$}] (a1) at (0.5, 3.4) {};
  \node[endpoint, fill=sourceblue!90!black, label={above right:\footnotesize $b_1$}] (b1) at (9.5, 3.4) {};

  % Interval 2
  \node[left=0.2cm, font=\small\bfseries, text=structurepurple!80!black] at (1.8, 2.3) {$I_2 = [a_2, b_2]$};
  \draw[intervalbar, draw=structurepurple!80!black, fill=structurepurple!15] (1.8, 2.1) rectangle (7.8, 2.5);
  \node[endpoint, fill=structurepurple!90!black, label={above left:\footnotesize $a_2$}] (a2) at (1.8, 2.3) {};
  \node[endpoint, fill=structurepurple!90!black, label={above right:\footnotesize $b_2$}] (b2) at (7.8, 2.3) {};

  % Interval 3
  \node[left=0.2cm, font=\small\bfseries, text=quotientorange!90!black] at (3.2, 1.2) {$I_3 = [a_3, b_3]$};
  \draw[intervalbar, draw=quotientorange!90!black, fill=quotientorange!18] (3.2, 1.0) rectangle (6.2, 1.4);
  \node[endpoint, fill=quotientorange!90!black, label={above left:\footnotesize $a_3$}] (a3) at (3.2, 1.2) {};
  \node[endpoint, fill=quotientorange!90!black, label={above right:\footnotesize $b_3$}] (b3) at (6.2, 1.2) {};

  % Limit point c
  \coordinate (c) at (4.7, 0.0);
  \node[endpoint, fill=targetgreen!90!black, inner sep=3pt] at (c) {};
  \node[below=3pt of c, font=\normalsize\bfseries, text=targetgreen!80!black] {$c = \sup_{n} a_n = \inf_{n} b_n$};

  % Convergence guide corridors
  \draw[guide, draw=sourceblue!50] (a1) -- (a2);
  \draw[guide, draw=structurepurple!50] (a2) -- (a3);
  \draw[guide, draw=quotientorange!70] (a3) -- (c);

  \draw[guide, draw=sourceblue!50] (b1) -- (b2);
  \draw[guide, draw=structurepurple!50] (b2) -- (b3);
  \draw[guide, draw=quotientorange!70] (b3) -- (c);

  % Central vertical target line
  \draw[densely dotted, thick, draw=targetgreen!80!black] (4.7, 3.8) -- (4.7, 0.2);
  \node[above=2pt, font=\footnotesize\bfseries, text=targetgreen!80!black] at (4.7, 3.8) {Intersection $\bigcap_n I_n = \{c\}$};
\end{tikzpicture}
\caption{The Nested Interval Theorem. A sequence of closed intervals $I_1 \supseteq I_2 \supseteq I_3 \supseteq \cdots$ shrinks so that $(b_n - a_n) \to 0$. The left endpoints $a_n$ are non-decreasing and bounded above by every right endpoint $b_m$. By the completeness of $\R$, the set of left endpoints has a least upper bound $c = \sup_n a_n = \inf_n b_n$, which is the unique point in the intersection $\bigcap_{n=1}^\infty I_n = \{c\}$. In $\Q$, the intersection of the corresponding rational intervals can be empty, as with rational approximations to $\sqrt{2}$.}\label{35:fig-nested-intervals}
\end{figure}

In the rational numbers $\Q$, this theorem fails completely. If we choose a sequence of nested intervals with rational endpoints trapping $\sqrt{2}$, say:
\[
[1, 2] \supset [1.4, 1.5] \supset [1.41, 1.42] \supset [1.414, 1.415] \supset \cdots
\]
their intersection inside $\Q$ is empty: $\bigcap_{n=1}^\infty (I_n \cap \Q) = \emptyset$. The rational numbers contain a gap where $\sqrt{2}$ should be, so the intersection in $\Q$ has no point. In contrast, the completeness of $\R$ ensures that every such nested sequence of closed intervals with vanishing length contains a unique real number.

\begin{recoverybox}
\textbf{Pause and think}: Why does the theorem require intervals to be \emph{closed} $[a_n, b_n]$? Consider the open intervals $J_n = (0, 1/n)$ on $\R$. They are nested: $(0, 1) \supset (0, 1/2) \supset (0, 1/3) \supset \cdots$, and their lengths shrink to zero. But what is $\bigcap_{n=1}^\infty J_n$? For any $x > 0$, by the Archimedean property there is some $n$ with $1/n \le x$, so $x \notin J_n$. And $0 \notin J_n$ by definition. The intersection is completely empty. The requirement that each interval be closed guarantees that the limit point belongs to every interval.
\end{recoverybox}

\subsection{Scratchpad versus Presentation: Working Backward, Writing Forward}

Writing an analysis proof requires switching between two opposing modes of thought:

\begin{enumerate}
  \item \textbf{The Scratchpad (Reverse Engineering)}: You begin by assuming the desired conclusion $|f(x) - L| < \eps$. You manipulate the expression backwards to see how small $|x - c|$ must be forced to become.
  \item \textbf{The Polished Proof (Forward Deduction)}: You discard the backwards scratchwork. You announce: ``Let $\eps > 0$ be given. Set $\delta = \dots$''. You then verify that whenever $0 < |x - c| < \delta$, the chain of inequalities leads forward to $|f(x) - L| < \eps$.
\end{enumerate}

Figure~\ref{35:fig-scratchpad-vs-proof} contrasts the two directions. The scratchpad works backwards from the target bound $|f(x) - L| < \eps$ to find what $\delta$ is needed. The polished proof runs forward: it states $\delta$ first, then proves that $|x - c| < \delta$ guarantees $|f(x) - L| < \eps$.

\begin{figure}[htbp]
\centering
\begin{tikzpicture}[scale=0.92, >=Latex]
  % PANEL 1: Scratchpad (Reverse Engineering)
  \begin{scope}[xshift=0cm]
    \node[anchor=south, font=\small\bfseries, text=quotientorange!90!black] at (2.7, 4.4) {1. Scratchpad (Backward Hunt)};
    \draw[rounded corners=6pt, fill=quotientorange!4, draw=quotientorange!50, thick] (0, -0.6) rectangle (5.4, 4.3);

    % Step 1: Goal
    \draw[rounded corners=3pt, fill=white, draw=quotientorange!70!black, thick] (0.3, 3.3) rectangle (5.1, 4.0);
    \node[font=\footnotesize\bfseries, text=quotientorange!95!black] at (2.7, 3.65) {Goal: $|f(x) - L| < \eps$};

    \draw[->, thick, quotientorange!80!black] (2.7, 3.3) -- (2.7, 2.9);

    % Step 2: Factor
    \draw[rounded corners=3pt, fill=white, draw=quotientorange!70!black, thick] (0.3, 2.2) rectangle (5.1, 2.9);
    \node[font=\footnotesize, text=gray!90!black] at (2.7, 2.55) {Factor: $|x - c| \cdot |g(x)| < \eps$};

    \draw[->, thick, quotientorange!80!black] (2.7, 2.2) -- (2.7, 1.8);

    % Step 3: Bound g(x)
    \draw[rounded corners=3pt, fill=white, draw=quotientorange!70!black, thick] (0.3, 1.1) rectangle (5.1, 1.8);
    \node[font=\footnotesize, text=gray!90!black] at (2.7, 1.45) {Bound: $\delta \le 1 \to |g(x)| \le M$};

    \draw[->, thick, quotientorange!80!black] (2.7, 1.1) -- (2.7, 0.7);

    % Step 4: Discovered Choice
    \draw[rounded corners=3pt, fill=quotientorange!15, draw=quotientorange!90!black, very thick] (0.3, 0.0) rectangle (5.1, 0.7);
    \node[font=\footnotesize\bfseries, text=quotientorange!95!black] at (2.7, 0.35) {Discover: $\delta = \min\{1, \, \eps/M\}$};

    \node[font=\scriptsize\bfseries, text=trapred!85!black] at (2.7, -0.3) {Reverse: Goal $\to$ Condition ($Q \to P$)};
  \end{scope}

  % PANEL 2: Polished Proof (Forward Deduction)
  \begin{scope}[xshift=6.6cm]
    \node[anchor=south, font=\small\bfseries, text=targetgreen!90!black] at (2.7, 4.4) {2. Polished Proof (Forward Deduction)};
    \draw[rounded corners=6pt, fill=targetgreen!4, draw=targetgreen!50, thick] (0, -0.6) rectangle (5.4, 4.3);

    % Step 1: Let eps > 0
    \draw[rounded corners=3pt, fill=white, draw=targetgreen!70!black, thick] (0.3, 3.3) rectangle (5.1, 4.0);
    \node[font=\footnotesize\bfseries, text=targetgreen!95!black] at (2.7, 3.65) {Given: Arbitrary $\eps > 0$};

    \draw[->, thick, targetgreen!80!black] (2.7, 3.3) -- (2.7, 2.9);

    % Step 2: Set delta
    \draw[rounded corners=3pt, fill=targetgreen!15, draw=targetgreen!90!black, very thick] (0.3, 2.2) rectangle (5.1, 2.9);
    \node[font=\footnotesize\bfseries, text=targetgreen!95!black] at (2.7, 2.55) {Announce: $\delta = \min\{1, \, \eps/M\}$};

    \draw[->, thick, targetgreen!80!black] (2.7, 2.2) -- (2.7, 1.8);

    % Step 3: Assume closeness
    \draw[rounded corners=3pt, fill=white, draw=targetgreen!70!black, thick] (0.3, 1.1) rectangle (5.1, 1.8);
    \node[font=\footnotesize, text=gray!90!black] at (2.7, 1.45) {Assume: $0 < |x - c| < \delta$};

    \draw[->, thick, targetgreen!80!black] (2.7, 1.1) -- (2.7, 0.7);

    % Step 4: Deduce conclusion
    \draw[rounded corners=3pt, fill=white, draw=targetgreen!70!black, thick] (0.3, 0.0) rectangle (5.1, 0.7);
    \node[font=\footnotesize\bfseries, text=targetgreen!95!black] at (2.7, 0.35) {Deduce: $|f(x) - L| < M\delta = \eps$};

    \node[font=\scriptsize\bfseries, text=targetgreen!85!black] at (2.7, -0.3) {Forward: Condition $\to$ Goal ($P \to Q$) \checkmark};
  \end{scope}
\end{tikzpicture}
\caption{The two opposing modes of analysis proof. \textbf{Left (The Scratchpad)}: You work backwards from the desired inequality to discover what $\delta$ must be. This is heuristic search, not a valid proof ($Q \implies P$). \textbf{Right (The Polished Proof)}: You discard the backwards scratchwork and argue strictly forward ($P \implies Q$), announcing $\delta$ up front and demonstrating that $|x - c| < \delta$ forces the output within the $\eps$ tolerance.}
\label{35:fig-scratchpad-vs-proof}
\end{figure}

To see this workflow in action, let us prove from first principles that $\lim_{x \to 2} (3x - 1) = 5$.

\paragraph{Stage 1: The Scratchpad.}
We want $|(3x - 1) - 5| < \eps$. Simplifying gives:
\[
|3x - 6| < \eps \iff 3|x - 2| < \eps \iff |x - 2| < \frac{\eps}{3}.
\]
This tells us that if we enforce $|x - 2| < \eps/3$, the desired bound will hold. We have found our candidate: $\delta = \eps/3$.

\paragraph{Stage 2: The Polished Forward Proof.}
Now we write down the forward argument, giving the reader no hint of the backwards exploration:
\begin{quote}
Let $\eps > 0$ be given. Choose $\delta = \eps/3$. Assume $0 < |x - 2| < \delta$. Then
\[
|(3x - 1) - 5| = |3x - 6| = 3|x - 2| < 3\delta = 3\left(\frac{\eps}{3}\right) = \eps.
\]
Because $\eps > 0$ was arbitrary, this proves that $\lim_{x \to 2} (3x - 1) = 5$.
\end{quote}

What happens when the function is nonlinear, like $f(x) = x^2$ near $c = 2$? On the scratchpad, we want:
\[
|x^2 - 4| = |x - 2| \cdot |x + 2| < \eps.
\]
Here we cannot simply divide by $|x + 2|$, because $\delta$ must be a fixed constant chosen before $x$ is known! We resolve this by putting a preliminary ceiling on $\delta$. We agree in advance that $\delta \le 1$. If $|x - 2| < 1$, then $1 < x < 3$, which guarantees that $|x + 2| < 5$. Under this constraint:
\[
|x^2 - 4| = |x - 2| \cdot |x + 2| < 5|x - 2|.
\]
To make $5|x - 2| < \eps$, we need $|x - 2| < \eps/5$. Combining our two requirements, we choose $\delta = \min\{1, \, \eps/5\}$. In the polished forward proof, you state this $\delta$ immediately and verify both bounds.

\begin{trapbox}
Never include your backwards scratch calculations inside the final proof as if they were deductions. If you write:
\[
|3x - 6| < \eps \implies 3|x - 2| < \eps \implies |x - 2| < \frac{\eps}{3},
\]
you have merely shown that \emph{if} the conclusion were true, \emph{then} $x$ would be close to $2$. A proof must show the exact reverse: that closeness of $x$ \emph{guarantees} the conclusion.
\end{trapbox}

\subsection{What Quantifier Discipline Unlocks---and Its Limits}

Weierstrassian quantifier discipline established rigorous foundations for key analysis results:
\begin{enumerate}
  \item \textbf{Legitimate Limit Interchange}: Quantifiers provide the exact theorem conditions under which derivatives, limits, and integrals commute: $\lim \int f_n = \int \lim f_n$ under uniform convergence.
  \item \textbf{The Core Theorems of Calculus}: The Intermediate Value Theorem, the Extreme Value Theorem, the Mean Value Theorem, and the Fundamental Theorem of Calculus all rest on Dedekind completeness and quantifier bounds.
  \item \textbf{Existence Theorems for Differential Equations}: The Picard--Lindel\"of theorem guarantees that ordinary differential equations modeling physical systems have unique, predictable trajectories. It proves this by converting differential equations into contraction mappings on complete spaces (the Banach Fixed-Point Theorem)---transformations that bring points closer together, converging to a unique fixed point.
\end{enumerate}

\paragraph{Scope and Limitations.}
Yet classical $\eps$-$\delta$ analysis has stark limitations:
\begin{itemize}
  \item \textbf{Algebraic Clutter and Metric Dependency}: $\eps$-$\delta$ proofs become unreadable in high dimensions or on curved geometric surfaces where Euclidean subtraction $|x - y|$ does not exist. The metric distance is often an arbitrary scaffolding that obscures the underlying qualitative behavior of space.
  \item \textbf{The Inadequacy of the Riemann Integral}: Even with uniform convergence, the classical Riemann integral fails for functions with dense discontinuities. For example, Dirichlet's indicator function of the rationals, $\mathbf{1}_\Q$ (which equals $1$ on rationals and $0$ on irrationals), oscillates too violently everywhere for rectangular Riemann approximations to settle. Pointwise limits of Riemann-integrable functions frequently fail to be integrable, preventing function spaces like $L^2$ (spaces of square-integrable wavefunctions and signals) from being complete.
\end{itemize}

\paragraph{Where We Go Next in the Book.}
Real analysis establishes properties of functions on $\R$ through precise inequalities and quantifiers. In Chapter~36 (\emph{Second Pass: Algebra as Lawful Manipulation}), we turn from continuous curves to symmetry and algebraic structure, studying operations governed by declared group laws rather than numerical approximations. In Chapter~37 (\emph{Second Pass: Topology as Stable Nearness}), we discard metrics entirely, studying continuity through open sets.

\subsection{Exercises}

\begin{enumerate}
  \item \textbf{Mental Check: The Quantifier Game}:
  \begin{enumerate}
    \item Negate the definition of uniform continuity: state precisely what an adversary must produce to prove that $f: D \to \R$ is \emph{not} uniformly continuous on $D$.
    \item Explain why uniform continuity on $D$ always implies pointwise continuity on $D$, but the reverse is false. Which quantifier moves where?
  \end{enumerate}
  \item \textbf{Uniform Continuity on an Unbounded Domain}:
  Let $f: [1, \infty) \to \R$ be given by $f(x) = 1/x$. Prove that $f$ is uniformly continuous on $[1, \infty)$. (Show that for all $x, y \ge 1$, $|f(x) - f(y)| \le |x - y|$, and explain why choosing $\delta = \eps$ works for every pair of points simultaneously.)
  \item \textbf{Scratchpad and Proof for a Quadratic Limit}:
  Prove from the definition of limit that $\lim_{x \to 3} (x^2 - 2x) = 3$. Provide your backwards scratchpad exploration first (including how you choose a preliminary ceiling on $\delta$ to bound $|x + 1|$), followed by the complete, forward deduction proof.
  \item \textbf{Testing the Boundaries of Nested Intervals}:
  Give explicit counterexamples showing why both hypotheses in the Nested Interval Theorem are necessary:
  \begin{enumerate}
    \item Find a sequence of nested open intervals $J_1 \supset J_2 \supset \cdots$ whose lengths shrink to zero, but whose intersection $\bigcap_{n=1}^\infty J_n$ is empty.
    \item Find a sequence of nested closed intervals $K_1 \supset K_2 \supset \cdots$ whose intersection is empty because the intervals are unbounded.
  \end{enumerate}
  \item \textbf{Cauchy Sequences and Completeness}:
  A sequence $(a_n)_{n=1}^\infty$ in $\R$ is a \emph{Cauchy sequence} if for every $\eps > 0$, there exists an index $N \in \N$ such that $|a_n - a_m| < \eps$ whenever $n, m \ge N$.
  \begin{enumerate}
    \item Prove that every convergent sequence in $\R$ is Cauchy.
    \item Give an example of a sequence of rational numbers that is Cauchy in $\Q$, but does not converge to any rational number in $\Q$. (This illustrates why completeness of $\R$ is equivalent to asserting that every Cauchy sequence converges.)
  \end{enumerate}
  \item \textbf{Uniform Continuity Preserves Cauchy Sequences}:
  Prove that if $f: \R \to \R$ is uniformly continuous and $(x_n)_{n=1}^\infty$ is a Cauchy sequence, then the image sequence $(f(x_n))_{n=1}^\infty$ is also a Cauchy sequence. Then show that this property fails for ordinary pointwise continuous functions by considering $f(x) = 1/x$ on the domain $(0, 1)$ with the sequence $x_n = 1/n$.
\end{enumerate}

\newpage
